\subsection{A complete high-ramification theorem}
\label{subsec:complete-high-ramification}

Throughout this subsection the ground field is algebraically closed of
characteristic zero.  Put
\[
 K_\epsilon=H_4+\epsilon H_3+\epsilon^2H_2+\epsilon^3H_1,
 \qquad
 D_j=[\epsilon^j]\det JK_\epsilon .
\]
For a Keller map, \(D_0=\cdots=D_8=0\) and
\(D_9=\det JH_1\ne0\).

Assume that the preceding reductions place the map in the primitive
coprime binary-pencil locus
\[
 H_4=(P,Q,0),\qquad P,Q\in k[x,y]_4,
 \qquad R=(H_3)_3\in k[x,y]_3,
\]
with
\[
 U=J(Q,R),\qquad V=J(P,R),\qquad W=J(P,Q)
\]
all nonzero.  Write
\[
 \Delta=\gcd(U,V,W),\qquad
 U=\Delta u,\quad V=\Delta v,\quad W=\Delta w,
 \qquad r=\deg\Delta.
\]
The planar minors give the Hilbert--Burch syzygy
\begin{equation}
 u\,dP-v\,dQ+w\,dR=0,
 \label{eq:r4-hb-syzygy}
\end{equation}
where
\[
 \deg u=\deg v=5-r,
 \qquad
 \deg w=6-r.
\]

\begin{theorem}[Primitive binary high ramification]
\label{thm:complete-binary-high-ramification}
In the preceding primitive coprime binary-pencil locus, if
\(r\ge4\), then \(F\) is a polynomial automorphism.
\end{theorem}

\begin{proof}
Since \(U,V\ne0\), one has \(r\le5\).  We treat \(r=5\) and
\(r=4\) separately.

\medskip
\noindent
\textbf{The case \(r=5\).}
Here \(u,v\) are constants and \(w\) is linear.  Set
\(S=uP-vQ\).  Equation \eqref{eq:r4-hb-syzygy} gives
\[
 dS=-w\,dR,
 \qquad
 dw\wedge dR=0.
\]
Homogeneity therefore yields
\[
 R=cw^3,
 \qquad
 S=-\frac{3c}{4}w^4.
 \tag{H.1}
\]
The full third coordinate has degree three and cubic part \(cw^3\).
The cubic-coordinate lemma straightens it while retaining \(w\) as one of
the two residual source coordinates.  The coordinate with leading part
\(-3cw^4/4\) then has residual plane degree at most nine: its degree-four
part is a polynomial in \(w\), while every term involving the eliminated
source variable has original degree at most three.  The two-variable
Appelgate--Onishi theorem, with its two key lemmas supplied by
Nowicki--Nakai, applies because \(9=3^2\).  Hence the residual plane map,
and therefore \(F\), is an automorphism.

\medskip
\noindent
\textbf{The case \(r=4\): dependent linear syzygies.}
Now \(u,v\) are linear and \(w\) is quadratic.  Suppose first that
\(u,v\) are dependent.  After changing the target basis of
\(\langle P,Q\rangle\), equation \eqref{eq:r4-hb-syzygy} becomes
\[
 \ell\,dS+w\,dR=0,
 \qquad
 4\ell S+3wR=0,                         \tag{H.2}
\]
with \(\ell\) linear.  If \(\ell\mid w\), write \(w=\ell m\).
Then \(dS=-m\,dR\), so \(dm\wedge dR=0\), and homogeneity again gives
\[
 R=cm^3,
 \qquad
 S=-\frac{3c}{4}m^4.
\]
This is the aligned situation (H.1).

If \(\ell\nmid w\), the Euler identity in (H.2) forces
\(R=\ell T\).  Substitution into (H.2) gives
\[
 d\!\left(\frac{\ell^4T}{w^3}\right)=0,
 \qquad
 \ell^4T=cw^3.
\]
Because \(\gcd(\ell,w)=1\), this is impossible in the unique
factorization domain \(k[x,y]\).  Thus the dependent branch is closed.

\medskip
\noindent
\textbf{The case \(r=4\): independent linear syzygies.}
Suppose now that \(u,v\) are independent.  A target change in
\(\langle P,Q\rangle\), followed by a source change in \((x,y)\), lets
us choose the two factors of the reduced quadratic \(w\) as the residual
linear syzygies.  If \(w\) is a square, the formulas below show that
\(P,Q\) have a common linear factor, contrary to primitivity.  Hence we
may normalize
\[
 u=x,
 \qquad
 v=y,
 \qquad
 w=xy.
\]
Euler contraction and differentiation of
\(x\,dP-y\,dQ+xy\,dR=0\) give
\begin{equation}
 P=-\frac14y^2R_y,
 \qquad
 Q=\frac14x^2R_x.                       \tag{H.3}
\end{equation}
Write
\[
 R=ax^3+bx^2y+cxy^2+dy^3.
\]
Coprimality of \(P,Q\) is equivalent to
\[
 a\ne0,
 \qquad
 d\ne0,
 \qquad
 \operatorname{Disc}(R)\ne0.           \tag{H.4}
\]
The common ramification quartic is
\begin{equation}
\begin{aligned}
2\Gamma={}&3abx^4+(9ac+b^2)x^3y+(18ad+4bc)x^2y^2\\
&\quad +(9bd+c^2)xy^3+3cdy^4,
\end{aligned}
\tag{H.5}
\end{equation}
and direct differentiation gives
\[
 J(Q,R)=x\Gamma,
 \qquad
 J(P,R)=y\Gamma,
 \qquad
 J(P,Q)=xy\Gamma.                       \tag{H.6}
\]

Write
\[
 H_3=(A,B,R),
 \qquad
 H_2=(C,D,E).
\]
The coefficient \(D_2=0\), together with (H.6), is
\[
 xA_z-yB_z+xyE_z=0.
\]
Consequently there are linear forms \(\ell,m\in k[x,y,z]_1\) such that
\[
 A_z=y(\ell-m/2),
 \qquad
 B_z=x(\ell+m/2),
 \qquad
 E_z=m.                                 \tag{H.7}
\]
Write \(\ell=\ell_0+\lambda z\) and \(m=m_0+\mu z\), with
\(\ell_0,m_0\in k[x,y]_1\).  The extreme coefficients of
\([z^3]D_3\) are
\[
 [y^0][z^3]D_3=\frac38a(2\lambda-\mu)^2x^3,
 \qquad
 [x^0][z^3]D_3=\frac38d(2\lambda+\mu)^2y^3.
\]
Condition (H.4) therefore gives \(\lambda=\mu=0\).  Thus
\(\ell,m\) are binary.

Put
\[
 \mathcal A=6R,
 \qquad
 \mathcal B=-3ax^3-bx^2y+cxy^2+3dy^3,
 \qquad
 \mathcal C=\frac12(3ax^3+bx^2y+cxy^2+3dy^3).
\]
One has the exact identity
\begin{equation}
 \mathcal A\mathcal C-\mathcal B^2=4xy\Gamma.       \tag{H.8}
\end{equation}
If \(u_2,v_2\) denote the coefficients of \(z^2\) in the first two
components of \(H_2\), then the coefficient of \(z\) in \(D_3\) is
\begin{equation}
 \mathcal Q_{\ell,m}+2\Gamma(u_2x-v_2y),
 \qquad
 \mathcal Q_{\ell,m}
 =\frac12(\mathcal A\ell^2+2\mathcal B\ell m+\mathcal C m^2).
 \tag{H.9}
\end{equation}
In particular,
\begin{equation}
 \Gamma\mid\mathcal Q_{\ell,m}.         \tag{H.10}
\end{equation}

\smallskip
\noindent
\emph{Squarefree \(\Gamma\).}
At a root of a squarefree \(\Gamma\), the symmetric matrix
\[
 M=\begin{pmatrix}\mathcal A&\mathcal B\\
                   \mathcal B&\mathcal C\end{pmatrix}
\]
has rank one.  It cannot vanish there: otherwise Euler's identity would
make that point a repeated root of \(R\), contrary to (H.4).  Hence
(H.10) implies that \((\ell,m)^T\) lies in the kernel of \(M\) at every
root of \(\Gamma\).  Since the two entries of
\(M(\ell,m)^T\) have degree four, there are constants \(\alpha,\beta\)
with
\[
 M\binom{\ell}{m}=\Gamma\binom{\alpha}{\beta}.
\]
Multiplication by the adjugate matrix and use of (H.8) give
\[
 4xy\ell=\mathcal C\alpha-\mathcal B\beta,
 \qquad
 4xym=-\mathcal B\alpha+\mathcal A\beta.
\]
Restriction to \(y=0\) gives \(\alpha+2\beta=0\), while restriction to
\(x=0\) gives \(\alpha-2\beta=0\).  Thus
\(\alpha=\beta=0\) and \(\ell=m=0\).

It remains to treat repeated \(\Gamma\).

\smallskip
\noindent
\emph{A repeated root away from \(xy=0\).}
A diagonal source change moves such a root to \([1:1]\).  The equations
\(\Gamma(1,1)=\partial_x\Gamma(1,1)=0\) have two components.  Their
eliminant is
\begin{equation}
 (b+2c+3d)
 (b^2+6bc+18bd+3c^2+12cd)=0.             \tag{H.11}
\end{equation}
On the first component, (H.4) forces
\[
 b=-2c-3d,
 \qquad
 a=c+2d,
\]
so \(R_x(1,1)=R_y(1,1)=0\), again contradicting primitivity.

The second component is the proper rational conic
\begin{equation}
\begin{aligned}
[a:b:c:d]=[&-(3u+2v)(3u^2+6uv+v^2):\\
&6u(3u+2v)(2u+3v):\\
&6v(3u+2v)(2u+3v):\\
&-(u^2+6uv+3v^2)(2u+3v)].
\end{aligned}
\tag{H.12}
\end{equation}
The four coordinates in (H.12) have no common projective zero, and the
inverse on the dense chart is \([u:v]=[b:c]\); hence (H.12) includes its
projective endpoints.  On the affine chart \(t=v/u\), put
\begin{equation}
\begin{aligned}
 a&=-(2t+3)(t^2+6t+3),\\
 b&=6(2t+3)(3t+2),\\
 c&=6t(2t+3)(3t+2),\\
 d&=-(3t+2)(3t^2+6t+1).
\end{aligned}
\tag{H.13}
\end{equation}
Then
\begin{equation}
 \Gamma=-9(2t+3)(3t+2)(x-y)^2G_{2,t},   \tag{H.14}
\end{equation}
where
\begin{equation}
\begin{aligned}
G_{2,t}={}&(2t^3+15t^2+24t+9)x^2\\
&+(6t^4+49t^3+90t^2+49t+6)xy\\
&+(9t^4+24t^3+15t^2+2t)y^2.
\end{aligned}
\tag{H.15}
\end{equation}
The relevant discriminants are
\begin{align}
 \operatorname{Disc}(R)
 &=3375(t+1)^6(2t+3)^2(3t+2)^2(3t^2+14t+3),
 \tag{H.16}\\
 \operatorname{Disc}(G_{2,t})
 &=(t+1)^2(2t+3)(3t+2)
   (2t^2+11t+2)(3t^2+14t+3),             \tag{H.17}\\
 G_{2,t}(1,1)
 &=15(t+1)^2(t^2+3t+1).                  \tag{H.18}
\end{align}
Thus, on the primitive locus, the only special residual divisors are the
\(3+1\) divisor \(t^2+3t+1=0\) and the internal \(2+2\) divisor
\(2t^2+11t+2=0\).  The projective point \(t=\infty\) is equivalent to
\(t=0\) under \(x\leftrightarrow y\), since
\[
 t^3(a,b,c,d)(1/t)=(d,c,b,a)(t).
\]

Away from the internal \(2+2\) divisor, exact row reduction of (H.10),
including the \(3+1\) specialization, gives the unique projective
first-normal direction
\begin{equation}
\begin{aligned}
 \ell&=h\bigl((t^2+6t+3)x-(3t^2+6t+1)y\bigr),\\
 m&=2h(t+1)\bigl((t+9)x+(9t+1)y\bigr).
\end{aligned}
\tag{H.19}
\end{equation}
Equation (H.9) then fixes
\[
 u_2=-\frac{2h^2(2t+3)}{3(3t+2)},
 \qquad
 v_2=\frac{2h^2t(3t+2)}{3(2t+3)}.
\]
All unrestricted binary terms and all remaining lower coefficients drop out
of the next highest normal coefficient, which factors as
\begin{equation}
\begin{aligned}
 [z^2]D_4={}&
 \frac{24h^3(t+1)}{(2t+3)(3t+2)}\\
 &\quad\cdot
 \bigl((2t+3)^2(4t+1)x+(t+4)(3t+2)^2y\bigr)G_{2,t}.
\end{aligned}
\tag{H.20}
\end{equation}
Every scalar and polynomial factor in (H.20) is nonzero on the primitive
locus.  Hence \(h=0\).

On the internal \(2+2\) divisor, work over
\[
 K=k[t]/(2t^2+11t+2).
\]
The complete first-normal space is two-dimensional:
\begin{equation}
\begin{aligned}
 L_0&=h,&
 L_1&=\frac{4t}{11}h-\frac3{22}k,\\
 M_0&=\frac{70}{11}h-
       \left(2+\frac{4t}{11}\right)k,&
 M_1&=k,
\end{aligned}
\tag{H.21}
\end{equation}
where \(\ell=L_0x+L_1y\) and \(m=M_0x+M_1y\).  After solving (H.9), the
extreme coefficients of \([z^2]D_4\) are
\begin{align}
 [x^3z^2]D_4
 &=\frac{12096}{1331}
 \left(h-\frac{2t+11}{24}k\right)^2
 \left(h+\frac{5(2t+11)}{56}k\right),       \tag{H.22}\\
 [y^3z^2]D_4
 &=-\frac{126(117t+22)}{1331}
 \left(h-\frac{2t+11}{2}k\right)^2
 \left(h+\frac{2t+11}{42}k\right).          \tag{H.23}
\end{align}
Both \(2t+11\) and \(117t+22\) are units in \(K\).  The projective root
sets
\[
 \left\{\frac1{24},-\frac5{56}\right\},
 \qquad
 \left\{\frac12,-\frac1{42}\right\}
\]
are disjoint, so (H.22)--(H.23) imply \(h=k=0\).

\smallskip
\noindent
\emph{A repeated root on \(xy=0\).}
If \(\Gamma\) has a repeated root at either endpoint, (H.4)--(H.5) force
\(b=c=0\).  After diagonal rescaling, take
\[
 R=x^3+y^3,
 \qquad
 \Gamma=9x^2y^2.
\]
The complete solution of (H.10) is
\[
 \ell=px+qy,
 \qquad
 m=2px-2qy,
 \qquad
 \frac{\mathcal Q_{\ell,m}}{\Gamma}
 =\frac43(q^2x+p^2y).
\]
After solving (H.9), the next coefficient is
\[
 [z^2]D_4=8q^3x^3-8p^3y^3,
\]
so \(p=q=0\).

We have now proved in every projective stratum that the first normal layer
vanishes.  Thus \(A,B,E\) are binary.  Returning to \(D_3=0\), the most
general remaining quadratic normal layer is
\[
 C=C_0(x,y)+\alpha yz,
 \qquad
 D=D_0(x,y)+\beta xz,
 \qquad
 (H_1)_3=L_{31}x+L_{32}y+(\beta-\alpha)z.
\]
The coefficient of \(z\) in \(D_5\), independently of every unrestricted
binary lower term, is
\begin{equation}
\begin{aligned}
\frac12\bigl(&6a\alpha^2x^3
 +(3b\alpha^2+2b\alpha\beta+b\beta^2)x^2y\\
&+(c\alpha^2+2c\alpha\beta+3c\beta^2)xy^2
 +6d\beta^2y^3\bigr).
\end{aligned}
\tag{H.24}
\end{equation}
The extreme coefficients and (H.4) give \(\alpha=\beta=0\).
Consequently all nonlinear terms of \(F\) are binary.

The third column of the invertible linear part is nonzero.  A target-linear
change sends it to \(e_3\), after which
\[
 F=(G_1(x,y),G_2(x,y),z+G_3(x,y)).
\]
The pair \((G_1,G_2)\) is a plane Keller map of degree at most four, so the
same Appelgate--Onishi--Nowicki--Nakai theorem makes it an automorphism.
This completes the proof.
\end{proof}

\begin{remark}[Evidence boundary]
All identities from (H.3) through (H.24), the projective conic
parametrization, the generic and \(3+1\) row reductions, the complete
\(2+2\) kernel, and the endpoint calculation are reconstructed by the
standalone exact checker
\begin{center}
\texttt{verify\_r4\_high\_ramification.py}.
\end{center}
It keeps arbitrary lower terms whenever they could enter the quoted
coefficient, uses exact rational arithmetic, makes no random
specializations, and divides only by factors explicitly excluded by
primitivity.  The structural Hilbert--Burch reduction into
\eqref{eq:r4-hb-syzygy} and the upstream routing into the primitive binary
locus remain separate inputs.
\end{remark}
